% =========================================================
% Chapter 4: Motion in Two Dimensions
% POSN 1 Physics Preparation Notes
% Language: Thai
% Compiler: XeLaTeX
% =========================================================

\chapter{การเคลื่อนที่ในสองมิติ}

\begin{center}
    {\Large \textbf{Motion in Two Dimensions}}\\[4pt]
    {\normalsize เอกสารเตรียมสอบคัดเลือก สอวน. ฟิสิกส์ ค่าย 1}\\[6pt]
    {\normalsize จัดทำโดย \textbf{Tames Thanakrit Damduan}}\\[2pt]
    {\footnotesize \textcopyright\ 2026 Tames Thanakrit Damduan. All rights reserved.}\\[8pt]
\end{center}

\section*{เป้าหมายของบทเรียน}

หลังจากเรียนบทนี้ นักเรียนควรทำสิ่งต่อไปนี้ได้

\begin{enumerate}
    \item เขียนตำแหน่ง ความเร็ว และความเร่งในรูปเวกเตอร์ได้
    \item หา $\vec{v}(t)$ และ $\vec{a}(t)$ จาก $\vec{r}(t)$ ได้
    \item แยกการเคลื่อนที่สองมิติเป็นแกน $x$ และแกน $y$ ได้
    \item ใช้สมการการเคลื่อนที่สองมิติเมื่อความเร่งคงตัวได้
    \item วิเคราะห์โพรเจคไทล์โดยแยกแนวราบและแนวดิ่งได้
    \item ใช้สูตรพื้นฐานของการเคลื่อนที่เป็นวงกลมสม่ำเสมอได้
    \item แยกความเร่งในการเคลื่อนที่ตามเส้นโค้งเป็นแนวสัมผัสและแนวรัศมีได้
    \item ใช้ความเร็วสัมพัทธ์ในการแก้โจทย์พื้นฐานได้
\end{enumerate}

\begin{tcolorbox}[title=แนวคิดสำคัญของบทนี้]
การเคลื่อนที่สองมิติไม่ได้เป็นเรื่องใหม่ทั้งหมด แต่เป็นการนำการเคลื่อนที่หนึ่งมิติมาคิดแยกตามแกน เมื่อแยกแกนถูกต้อง แกน $x$ และแกน $y$ จะใช้สมการเดิมจากบทการเคลื่อนที่แนวตรง โดยมีเวลาเดียวกันเป็นตัวเชื่อม
\end{tcolorbox}

\newpage{}

% =========================================================
\section{เวกเตอร์ตำแหน่ง ความเร็ว และความเร่ง}

\subsection{เวกเตอร์ตำแหน่ง}

ในหนึ่งมิติ เราบอกตำแหน่งด้วยฟังก์ชัน

\[
x(t)
\]

แต่ในสองมิติ ต้องบอกทั้งตำแหน่งตามแกน $x$ และตำแหน่งตามแกน $y$ จึงเขียนตำแหน่งเป็นเวกเตอร์

\[
\vec{r}(t)=x(t)\hat{i}+y(t)\hat{j}
\]

โดย

\[
x(t)=\text{ตำแหน่งตามแกน }x
\]

\[
y(t)=\text{ตำแหน่งตามแกน }y
\]

\begin{tcolorbox}[title=นิยาม]
เวกเตอร์ตำแหน่ง $\vec{r}(t)$ คือเวกเตอร์ที่ลากจากจุดกำเนิดไปยังตำแหน่งของวัตถุ ณ เวลา $t$
\end{tcolorbox}

\subsection{เวกเตอร์การกระจัด}

ถ้าวัตถุเปลี่ยนตำแหน่งจาก

\[
\vec{r}_i
\]

ไปยัง

\[
\vec{r}_f
\]

การกระจัดคือ

\[
\Delta \vec{r}=\vec{r}_f-\vec{r}_i
\]

ถ้า

\[
\vec{r}_i=x_i\hat{i}+y_i\hat{j}
\]

และ

\[
\vec{r}_f=x_f\hat{i}+y_f\hat{j}
\]

จะได้

\[
\Delta \vec{r}
=
(x_f-x_i)\hat{i}+(y_f-y_i)\hat{j}
\]

ดังนั้น

\[
\Delta \vec{r}=\Delta x\hat{i}+\Delta y\hat{j}
\]

\begin{tcolorbox}[title=สรุป]
\[
\Delta \vec{r}=\vec{r}_f-\vec{r}_i
\]

\[
\Delta \vec{r}=\Delta x\hat{i}+\Delta y\hat{j}
\]
\end{tcolorbox}

\subsection{ความเร็วเฉลี่ย}

ในหนึ่งมิติ

\[
v_{\text{avg}}=\frac{\Delta x}{\Delta t}
\]

ในสองมิติ ให้เปลี่ยนจากการกระจัดหนึ่งมิติเป็นเวกเตอร์การกระจัด

\[
\vec{v}_{\text{avg}}=\frac{\Delta \vec{r}}{\Delta t}
\]

แทนรูปองค์ประกอบ

\[
\vec{v}_{\text{avg}}
=
\frac{\Delta x}{\Delta t}\hat{i}
+
\frac{\Delta y}{\Delta t}\hat{j}
\]

ดังนั้น

\[
\vec{v}_{\text{avg}}
=
v_{x,\text{avg}}\hat{i}
+
v_{y,\text{avg}}\hat{j}
\]

\begin{tcolorbox}[title=ข้อควรจำ]
ความเร็วเฉลี่ยมีทิศเดียวกับการกระจัด ไม่จำเป็นต้องมีทิศเดียวกับเส้นทางจริงที่วัตถุเคลื่อนที่
\end{tcolorbox}

\subsection{ความเร็วขณะหนึ่ง}

ความเร็วขณะหนึ่งคือค่าที่ได้เมื่อทำช่วงเวลาให้เล็กลงจนเข้าใกล้ศูนย์

\[
\vec{v}
=
\lim_{\Delta t\to 0}
\frac{\Delta \vec{r}}{\Delta t}
\]

ในแคลคูลัสเขียนเป็น

\[
\vec{v}(t)=\frac{d\vec{r}}{dt}
\]

ถ้า

\[
\vec{r}(t)=x(t)\hat{i}+y(t)\hat{j}
\]

และแกนไม่หมุนตามเวลา จะหาอนุพันธ์ทีละแกนได้

\[
\vec{v}(t)=\frac{dx}{dt}\hat{i}+\frac{dy}{dt}\hat{j}
\]

หรือ

\[
\vec{v}(t)=v_x(t)\hat{i}+v_y(t)\hat{j}
\]

\begin{tcolorbox}[title=ข้อควรจำ]
ความเร็วขณะหนึ่งมีทิศสัมผัสกับเส้นทางการเคลื่อนที่ ณ จุดนั้นเสมอ
\end{tcolorbox}

\subsection{ความเร่ง}

ความเร่งเฉลี่ยในสองมิติคือ

\[
\vec{a}_{\text{avg}}=\frac{\Delta \vec{v}}{\Delta t}
\]

โดย

\[
\Delta \vec{v}=\vec{v}_f-\vec{v}_i
\]

ความเร่งขณะหนึ่งคือ

\[
\vec{a}
=
\lim_{\Delta t\to 0}
\frac{\Delta \vec{v}}{\Delta t}
\]

ดังนั้น

\[
\vec{a}(t)=\frac{d\vec{v}}{dt}
\]

และเพราะ

\[
\vec{v}(t)=\frac{d\vec{r}}{dt}
\]

จึงได้

\[
\vec{a}(t)=\frac{d^2\vec{r}}{dt^2}
\]

ถ้า

\[
\vec{v}(t)=v_x(t)\hat{i}+v_y(t)\hat{j}
\]

จะได้

\[
\vec{a}(t)=\frac{dv_x}{dt}\hat{i}+\frac{dv_y}{dt}\hat{j}
\]

หรือ

\[
\vec{a}(t)=a_x(t)\hat{i}+a_y(t)\hat{j}
\]

\begin{examplebox}{ตัวอย่าง: หา $\vec{v}(t)$ และ $\vec{a}(t)$ จาก $\vec{r}(t)$}
กำหนดให้

\[
\vec{r}(t)=(t^2+2t)\hat{i}+(4t-t^2)\hat{j}
\]

หาความเร็วจาก

\[
\vec{v}(t)=\frac{d\vec{r}}{dt}
\]

จึงได้

\[
\vec{v}(t)=(2t+2)\hat{i}+(4-2t)\hat{j}
\]

หาความเร่งจาก

\[
\vec{a}(t)=\frac{d\vec{v}}{dt}
\]

จึงได้

\[
\vec{a}(t)=2\hat{i}-2\hat{j}
\]
\end{examplebox}

\begin{exercisebox}{แบบฝึกหัด 4.1}
\begin{enumerate}
    \item อนุภาคมี $\vec{r}_i=2\hat{i}+3\hat{j}$ และ $\vec{r}_f=8\hat{i}-1\hat{j}$ จงหา $\Delta \vec{r}$
    \item ถ้าอนุภาคในข้อก่อนหน้าใช้เวลา $2\,\si{s}$ จงหา $\vec{v}_{\text{avg}}$
    \item ถ้า $\vec{r}(t)=5t\hat{i}+t^2\hat{j}$ จงหา $\vec{v}(t)$ และ $\vec{a}(t)$
    \item ถ้า $\vec{v}(t)=3t^2\hat{i}-2t\hat{j}$ จงหา $\vec{a}(t)$
    \item ถ้า $\vec{r}(t)=(t^3-2t)\hat{i}+5t^2\hat{j}$ จงหา $\vec{v}(1)$
\end{enumerate}
\end{exercisebox}

% =========================================================
\section{การเคลื่อนที่สองมิติเมื่อความเร่งคงตัว}

\subsection{สมการเวกเตอร์}

ในหนึ่งมิติ เมื่อความเร่งคงตัว เรามี

\[
v_f=v_i+at
\]

และ

\[
x_f=x_i+v_it+\frac{1}{2}at^2
\]

ในสองมิติ สมการมีรูปเหมือนเดิม แต่เปลี่ยนเป็นเวกเตอร์

\[
\vec{v}_f=\vec{v}_i+\vec{a}t
\]

\[
\vec{r}_f=\vec{r}_i+\vec{v}_it+\frac{1}{2}\vec{a}t^2
\]

\begin{tcolorbox}[title=สมการเวกเตอร์เมื่อความเร่งคงตัว]
\[
\vec{v}_f=\vec{v}_i+\vec{a}t
\]

\[
\vec{r}_f=\vec{r}_i+\vec{v}_it+\frac{1}{2}\vec{a}t^2
\]
\end{tcolorbox}

\subsection{แยกเป็นสมการตามแกน}

ให้

\[
\vec{v}_i=v_{xi}\hat{i}+v_{yi}\hat{j}
\]

และ

\[
\vec{a}=a_x\hat{i}+a_y\hat{j}
\]

จากสมการความเร็ว

\[
\vec{v}_f=\vec{v}_i+\vec{a}t
\]

จะได้

\[
v_{xf}=v_{xi}+a_xt
\]

\[
v_{yf}=v_{yi}+a_yt
\]

จากสมการตำแหน่ง

\[
\vec{r}_f=\vec{r}_i+\vec{v}_it+\frac{1}{2}\vec{a}t^2
\]

จะได้

\[
x_f=x_i+v_{xi}t+\frac{1}{2}a_xt^2
\]

\[
y_f=y_i+v_{yi}t+\frac{1}{2}a_yt^2
\]

\begin{tcolorbox}[title=สมการตามแกนเมื่อความเร่งคงตัว]
แกน $x$

\[
v_{xf}=v_{xi}+a_xt
\]

\[
x_f=x_i+v_{xi}t+\frac{1}{2}a_xt^2
\]

แกน $y$

\[
v_{yf}=v_{yi}+a_yt
\]

\[
y_f=y_i+v_{yi}t+\frac{1}{2}a_yt^2
\]
\end{tcolorbox}

\subsection{หลักคิดสำคัญ}

แกน $x$ และแกน $y$ คิดแยกกันได้ แต่ต้องใช้เวลาเดียวกัน

\[
t_x=t_y
\]

ในโจทย์จำนวนมาก เราจะใช้แกนหนึ่งหาเวลา แล้วนำเวลานั้นไปใช้กับอีกแกนหนึ่ง

\begin{examplebox}{ตัวอย่าง: ความเร่งคงตัวในสองมิติ}
อนุภาคเริ่มจากจุดกำเนิดด้วยความเร็วต้น

\[
\vec{v}_i=4\hat{i}+6\hat{j}
\]

และมีความเร่งคงตัว

\[
\vec{a}=2\hat{i}-3\hat{j}
\]

จงหาตำแหน่งเมื่อ

\[
t=2\,\si{s}
\]

ใช้สมการ

\[
\vec{r}_f=\vec{r}_i+\vec{v}_it+\frac{1}{2}\vec{a}t^2
\]

เพราะเริ่มจากจุดกำเนิด

\[
\vec{r}_i=\vec{0}
\]

แทนค่า

\[
\vec{r}_f=(4\hat{i}+6\hat{j})(2)+\frac{1}{2}(2\hat{i}-3\hat{j})(2^2)
\]

\[
\vec{r}_f=8\hat{i}+12\hat{j}+4\hat{i}-6\hat{j}
\]

ดังนั้น

\[
\vec{r}_f=12\hat{i}+6\hat{j}
\]
\end{examplebox}

\begin{exercisebox}{แบบฝึกหัด 4.2}
\begin{enumerate}
    \item อนุภาคมี $\vec{v}_i=3\hat{i}+4\hat{j}$ และ $\vec{a}=2\hat{i}-1\hat{j}$ จงหา $\vec{v}_f$ เมื่อ $t=5\,\si{s}$
    \item อนุภาคเริ่มที่จุดกำเนิด มี $\vec{v}_i=2\hat{i}+6\hat{j}$ และ $\vec{a}=0\hat{i}-2\hat{j}$ จงหา $\vec{r}(3)$
    \item ถ้า $a_x=0$ และ $a_y=-g$ จงอธิบายว่าการเคลื่อนที่นี้ตรงกับสถานการณ์ใด
    \item ถ้า $x=4t$ และ $y=3t-5t^2$ จงกำจัด $t$ เพื่อเขียน $y$ เป็นฟังก์ชันของ $x$
\end{enumerate}
\end{exercisebox}

% =========================================================
\section{โพรเจคไทล์}

\subsection{แนวคิดหลัก}

โพรเจคไทล์คือวัตถุที่ถูกยิงหรือขว้าง แล้วหลังจากนั้นมีความเร่งจากแรงโน้มถ่วงเพียงอย่างเดียว โดยละเลยแรงต้านอากาศ

ถ้าเลือกแกน $x$ เป็นแนวราบไปทางขวา และแกน $y$ เป็นแนวดิ่งขึ้น จะได้

\[
a_x=0
\]

\[
a_y=-g
\]

ดังนั้น

\[
\text{แกน }x:\text{ ความเร็วคงตัว}
\]

\[
\text{แกน }y:\text{ ความเร่งคงตัว } -g
\]

\begin{tcolorbox}[title=หัวใจของโพรเจคไทล์]
แกน $x$ และแกน $y$ คิดแยกกัน แต่ใช้เวลาเดียวกัน
\end{tcolorbox}

\subsection{แตกความเร็วต้น}

ถ้าวัตถุถูกยิงด้วยความเร็วต้นขนาด $u$ ทำมุม $\theta$ กับแนวระดับ จะมี

\[
u_x=u\cos\theta
\]

\[
u_y=u\sin\theta
\]

ดังนั้นเวกเตอร์ความเร็วต้นคือ

\[
\vec{u}=u\cos\theta\hat{i}+u\sin\theta\hat{j}
\]

\begin{tcolorbox}[title=องค์ประกอบความเร็วต้น]
\[
u_x=u\cos\theta
\]

\[
u_y=u\sin\theta
\]
\end{tcolorbox}

\subsection{สมการโพรเจคไทล์จากจุดกำเนิด}

ให้จุดเริ่มต้นเป็นจุดกำเนิด

\[
x_i=0,\quad y_i=0
\]

แกน $x$

\[
x=u_xt
\]

แทน

\[
u_x=u\cos\theta
\]

จะได้

\[
x=u\cos\theta\,t
\]

แกน $y$

\[
y=u_yt-\frac{1}{2}gt^2
\]

แทน

\[
u_y=u\sin\theta
\]

จะได้

\[
y=u\sin\theta\,t-\frac{1}{2}gt^2
\]

ความเร็วแต่ละแกนคือ

\[
v_x=u\cos\theta
\]

\[
v_y=u\sin\theta-gt
\]

\begin{tcolorbox}[title=สมการโพรเจคไทล์]
\[
x=u\cos\theta\,t
\]

\[
y=u\sin\theta\,t-\frac{1}{2}gt^2
\]

\[
v_x=u\cos\theta
\]

\[
v_y=u\sin\theta-gt
\]
\end{tcolorbox}

\subsection{สมการเส้นทาง}

เริ่มจาก

\[
x=u\cos\theta\,t
\]

จัดรูปหาเวลา

\[
t=\frac{x}{u\cos\theta}
\]

แทนใน

\[
y=u\sin\theta\,t-\frac{1}{2}gt^2
\]

จะได้

\[
y=u\sin\theta\left(\frac{x}{u\cos\theta}\right)
-\frac{1}{2}g\left(\frac{x}{u\cos\theta}\right)^2
\]

ดังนั้น

\[
y=x\tan\theta-\frac{gx^2}{2u^2\cos^2\theta}
\]

สมการนี้มีพจน์ $x^2$ จึงเป็นพาราโบลา

\begin{tcolorbox}[title=สมการเส้นทางของโพรเจคไทล์]
\[
y=x\tan\theta-\frac{gx^2}{2u^2\cos^2\theta}
\]
\end{tcolorbox}

\subsection{จุดสูงสุด}

ที่จุดสูงสุด ความเร็วแนวดิ่งเป็นศูนย์

\[
v_y=0
\]

จาก

\[
v_y=u\sin\theta-gt
\]

จะได้เวลาถึงจุดสูงสุด

\[
t_{\text{top}}=\frac{u\sin\theta}{g}
\]

ความสูงสูงสุดหาได้จาก

\[
v_y^2=u_y^2+2a_y\Delta y
\]

ที่จุดสูงสุด

\[
v_y=0,\quad u_y=u\sin\theta,\quad a_y=-g
\]

ดังนั้น

\[
0=(u\sin\theta)^2-2gH
\]

จึงได้

\[
H=\frac{u^2\sin^2\theta}{2g}
\]

\begin{tcolorbox}[title=จุดสูงสุด]
\[
t_{\text{top}}=\frac{u\sin\theta}{g}
\]

\[
H=\frac{u^2\sin^2\theta}{2g}
\]
\end{tcolorbox}

\subsection{กรณียิงและตกที่ระดับเดิม}

ถ้าวัตถุเริ่มและตกที่ระดับเดียวกัน จะมี

\[
y=0
\]

ใช้

\[
y=u\sin\theta\,t-\frac{1}{2}gt^2
\]

แทน

\[
0=u\sin\theta\,t-\frac{1}{2}gt^2
\]

ดึง $t$ เป็นตัวร่วม

\[
0=t\left(u\sin\theta-\frac{1}{2}gt\right)
\]

ได้คำตอบ

\[
t=0
\]

ซึ่งเป็นจุดเริ่มต้น และเวลาทั้งหมดคือ

\[
T=\frac{2u\sin\theta}{g}
\]

พิสัยแนวราบคือ

\[
R=u\cos\theta\,T
\]

แทนค่า $T$

\[
R=u\cos\theta\left(\frac{2u\sin\theta}{g}\right)
\]

\[
R=\frac{u^2\sin2\theta}{g}
\]

\begin{tcolorbox}[title=กรณียิงและตกที่ระดับเดิม]
\[
T=\frac{2u\sin\theta}{g}
\]

\[
R=\frac{u^2\sin2\theta}{g}
\]
\end{tcolorbox}

\subsection{ข้อควรระวัง}

สูตร

\[
R=\frac{u^2\sin2\theta}{g}
\]

ใช้ได้เฉพาะกรณีที่จุดเริ่มต้นและจุดตกอยู่ระดับเดียวกันเท่านั้น ถ้าตกบนพื้นเอียง ตกจากหน้าผา หรือจุดตกสูงต่างจากจุดเริ่มต้น ต้องกลับไปใช้สมการพื้นฐาน

\[
x=u\cos\theta\,t
\]

\[
y=u\sin\theta\,t-\frac{1}{2}gt^2
\]

\begin{examplebox}{ตัวอย่าง: โพรเจคไทล์พื้นฐาน}
ยิงวัตถุด้วยความเร็วต้น $20\,\si{m/s}$ ทำมุม $30^\circ$ กับแนวระดับ ใช้ $g=10\,\si{m/s^2}$ จงหาเวลาถึงจุดสูงสุดและความสูงสูงสุด

องค์ประกอบแนวดิ่งคือ

\[
u_y=20\sin30^\circ
\]

\[
u_y=10\,\si{m/s}
\]

เวลาถึงจุดสูงสุดคือ

\[
t_{\text{top}}=\frac{u_y}{g}
\]

\[
t_{\text{top}}=\frac{10}{10}=1\,\si{s}
\]

ความสูงสูงสุดคือ

\[
H=\frac{u_y^2}{2g}
\]

\[
H=\frac{10^2}{2(10)}=5\,\si{m}
\]
\end{examplebox}

\begin{examplebox}{ตัวอย่าง: ยิงในแนวระดับจากหน้าผา}
วัตถุถูกยิงในแนวระดับจากหน้าผาสูง $45\,\si{m}$ ด้วยความเร็ว $10\,\si{m/s}$ ใช้ $g=10\,\si{m/s^2}$ จงหาระยะทางแนวราบก่อนตกถึงพื้น

แนวดิ่งเริ่มจาก

\[
u_y=0
\]

และตกลงเป็นระยะ

\[
45\,\si{m}
\]

ใช้

\[
45=\frac{1}{2}gt^2
\]

แทนค่า

\[
45=\frac{1}{2}(10)t^2
\]

\[
45=5t^2
\]

\[
t=3\,\si{s}
\]

แนวราบมีความเร็วคงตัว

\[
x=u_xt
\]

\[
x=(10)(3)=30\,\si{m}
\]
\end{examplebox}

\begin{exercisebox}{แบบฝึกหัด 4.3}
\begin{enumerate}
    \item ยิงวัตถุด้วย $u=20\,\si{m/s}$ ทำมุม $45^\circ$ กับแนวระดับ ใช้ $g=10\,\si{m/s^2}$ จงหา $u_x$ และ $u_y$
    \item จากข้อก่อนหน้า จงหาเวลาทั้งหมดถ้าวัตถุตกที่ระดับเดิม
    \item วัตถุถูกยิงในแนวระดับจากหน้าผาสูง $45\,\si{m}$ ด้วยความเร็ว $10\,\si{m/s}$ ใช้ $g=10\,\si{m/s^2}$ จงหาระยะทางแนวราบก่อนตกถึงพื้น
    \item จงพิสูจน์ว่าสำหรับการยิงและตกที่ระดับเดิม มุม $45^\circ$ ให้พิสัยมากที่สุด
    \item โพรเจคไทล์มี $u_x=12\,\si{m/s}$ และ $u_y=16\,\si{m/s}$ ใช้ $g=10\,\si{m/s^2}$ จงหาความสูงสูงสุด
\end{enumerate}
\end{exercisebox}

% =========================================================
\section{การเคลื่อนที่เป็นวงกลมสม่ำเสมอ}

\subsection{อัตราเร็วคงตัวแต่ความเร่งไม่เป็นศูนย์}

ในการเคลื่อนที่เป็นวงกลมสม่ำเสมอ วัตถุมีอัตราเร็วคงตัว

\[
v=\text{constant}
\]

แต่ทิศของเวกเตอร์ความเร็วเปลี่ยนตลอดเวลา เพราะความเร็วต้องสัมผัสกับเส้นทางวงกลม

ดังนั้นแม้ขนาดของความเร็วไม่เปลี่ยน แต่เวกเตอร์ความเร็วเปลี่ยน จึงมีความเร่ง

\begin{tcolorbox}[title=จุดสำคัญ]
อัตราเร็วคงตัวไม่ได้แปลว่าความเร่งเป็นศูนย์ ถ้าทิศของความเร็วเปลี่ยน แสดงว่ามีความเร่ง
\end{tcolorbox}

\subsection{ความเร่งสู่ศูนย์กลาง}

ความเร่งของการเคลื่อนที่เป็นวงกลมสม่ำเสมอชี้เข้าหาศูนย์กลางวงกลมเสมอ เรียกว่า ความเร่งสู่ศูนย์กลาง

\[
\vec{a}_c \text{ ชี้เข้าสู่ศูนย์กลาง}
\]

ขนาดของความเร่งสู่ศูนย์กลางคือ

\[
a_c=\frac{v^2}{r}
\]

โดยที่ $r$ คือรัศมีของวงกลม

\begin{tcolorbox}[title=ความเร่งสู่ศูนย์กลาง]
\[
a_c=\frac{v^2}{r}
\]

ทิศของ $\vec{a}_c$ ชี้เข้าสู่ศูนย์กลางวงกลม
\end{tcolorbox}

\subsection{คาบและอัตราเร็วเชิงมุม}

คาบ $T$ คือเวลาที่วัตถุใช้เคลื่อนที่ครบหนึ่งรอบ

ระยะทางหนึ่งรอบคือ

\[
2\pi r
\]

ดังนั้นอัตราเร็วคือ

\[
v=\frac{2\pi r}{T}
\]

อัตราเร็วเชิงมุมคือมุมที่กวาดได้ต่อเวลา

\[
\omega=\frac{2\pi}{T}
\]

เนื่องจาก

\[
v=\frac{2\pi r}{T}
\]

และ

\[
\omega=\frac{2\pi}{T}
\]

จึงได้

\[
v=\omega r
\]

แทนใน

\[
a_c=\frac{v^2}{r}
\]

จะได้

\[
a_c=\omega^2r
\]

\begin{tcolorbox}[title=สูตรวงกลมสม่ำเสมอ]
\[
v=\frac{2\pi r}{T}
\]

\[
\omega=\frac{2\pi}{T}
\]

\[
v=\omega r
\]

\[
a_c=\frac{v^2}{r}=\omega^2r
\]
\end{tcolorbox}

\begin{examplebox}{ตัวอย่าง: วงกลมสม่ำเสมอ}
วัตถุเคลื่อนที่เป็นวงกลมรัศมี $2\,\si{m}$ ด้วยอัตราเร็ว $6\,\si{m/s}$ จงหาความเร่งสู่ศูนย์กลาง

ใช้

\[
a_c=\frac{v^2}{r}
\]

แทนค่า

\[
a_c=\frac{6^2}{2}
\]

\[
a_c=18\,\si{m/s^2}
\]

ทิศของความเร่งชี้เข้าสู่ศูนย์กลางวงกลม
\end{examplebox}

\begin{exercisebox}{แบบฝึกหัด 4.4}
\begin{enumerate}
    \item วัตถุเคลื่อนที่เป็นวงกลมรัศมี $5\,\si{m}$ ด้วยอัตราเร็ว $10\,\si{m/s}$ จงหา $a_c$
    \item วัตถุเคลื่อนที่ครบหนึ่งรอบในเวลา $4\,\si{s}$ บนวงกลมรัศมี $2\,\si{m}$ จงหา $v$
    \item ถ้า $\omega=3\,\si{rad/s}$ และ $r=4\,\si{m}$ จงหา $v$ และ $a_c$
    \item ถ้าวัตถุมีอัตราเร็วเพิ่มเป็นสองเท่า แต่รัศมีเท่าเดิม ความเร่งสู่ศูนย์กลางเปลี่ยนเป็นกี่เท่า
\end{enumerate}
\end{exercisebox}

% =========================================================
\section{ความเร่งแนวสัมผัสและแนวรัศมี}

\subsection{การเคลื่อนที่ตามเส้นโค้งที่อัตราเร็วเปลี่ยน}

ถ้าวัตถุเคลื่อนที่ตามเส้นโค้งและอัตราเร็วเปลี่ยนด้วย ความเร่งจะมีสองหน้าที่

\[
\text{เปลี่ยนขนาดของความเร็ว}
\]

และ

\[
\text{เปลี่ยนทิศของความเร็ว}
\]

จึงแยกความเร่งเป็นสองส่วน

\[
\vec{a}=\vec{a}_t+\vec{a}_r
\]

โดย

\[
\vec{a}_t=\text{ความเร่งแนวสัมผัส}
\]

\[
\vec{a}_r=\text{ความเร่งแนวรัศมี}
\]

\subsection{ความเร่งแนวสัมผัส}

ความเร่งแนวสัมผัสทำหน้าที่เปลี่ยนขนาดของความเร็ว

\[
a_t=\left|\frac{dv}{dt}\right|
\]

ถ้าอัตราเร็วเพิ่มขึ้น $\vec{a}_t$ มีทิศเดียวกับการเคลื่อนที่

ถ้าอัตราเร็วลดลง $\vec{a}_t$ มีทิศตรงข้ามกับการเคลื่อนที่

\subsection{ความเร่งแนวรัศมี}

ความเร่งแนวรัศมีทำหน้าที่เปลี่ยนทิศของความเร็ว

\[
a_r=\frac{v^2}{r}
\]

ทิศของ $\vec{a}_r$ ชี้เข้าหาศูนย์กลางความโค้ง

\begin{tcolorbox}[title=ความเร่งในการเคลื่อนที่ตามเส้นโค้ง]
\[
a_t=\left|\frac{dv}{dt}\right|
\]

\[
a_r=\frac{v^2}{r}
\]

ถ้า $\vec{a}_t$ และ $\vec{a}_r$ ตั้งฉากกัน

\[
a=\sqrt{a_t^2+a_r^2}
\]
\end{tcolorbox}

\begin{examplebox}{ตัวอย่าง: ความเร่งสองส่วน}
รถเคลื่อนที่บนทางโค้งรัศมี $50\,\si{m}$ ด้วยอัตราเร็ว $20\,\si{m/s}$ และกำลังเพิ่มอัตราเร็วด้วยอัตรา $2\,\si{m/s^2}$ จงหาขนาดความเร่งรวม

ความเร่งแนวสัมผัสคือ

\[
a_t=2\,\si{m/s^2}
\]

ความเร่งแนวรัศมีคือ

\[
a_r=\frac{v^2}{r}
\]

\[
a_r=\frac{20^2}{50}=8\,\si{m/s^2}
\]

เพราะสององค์ประกอบตั้งฉากกัน

\[
a=\sqrt{a_t^2+a_r^2}
\]

\[
a=\sqrt{2^2+8^2}
\]

\[
a=2\sqrt{17}\,\si{m/s^2}
\]
\end{examplebox}

\begin{exercisebox}{แบบฝึกหัด 4.5}
\begin{enumerate}
    \item วัตถุเคลื่อนที่บนทางโค้งรัศมี $10\,\si{m}$ ด้วยอัตราเร็ว $5\,\si{m/s}$ และอัตราเร็วคงตัว จงหา $a_t$, $a_r$, และ $a$
    \item รถเคลื่อนที่บนโค้งรัศมี $80\,\si{m}$ ด้วยอัตราเร็ว $20\,\si{m/s}$ และกำลังชะลอด้วย $3\,\si{m/s^2}$ จงหาขนาดความเร่งรวม
    \item ถ้าวัตถุเคลื่อนที่เป็นเส้นตรงแต่อัตราเร็วเพิ่มขึ้น จะมี $a_r$ หรือไม่ เพราะเหตุใด
    \item ถ้าวัตถุเคลื่อนที่เป็นวงกลมด้วยอัตราเร็วคงตัว จะมี $a_t$ หรือไม่ เพราะเหตุใด
\end{enumerate}
\end{exercisebox}

% =========================================================
\section{ความเร็วสัมพัทธ์และความเร่งสัมพัทธ์}

\subsection{แนวคิดหลัก}

ผู้สังเกตที่อยู่ในกรอบอ้างอิงต่างกันอาจวัดความเร็วของวัตถุเดียวกันได้ต่างกัน

ตัวอย่างเช่น คนบนเรืออาจเห็นวัตถุลอยน้ำอยู่ใกล้ตัว แต่คนบนฝั่งเห็นทั้งเรือและวัตถุเคลื่อนที่ไปพร้อมกับกระแสน้ำ

ดังนั้นต้องระบุให้ชัดว่า

\[
\text{ความเร็วของอะไร เทียบกับอะไร}
\]

\subsection{ตำแหน่งสัมพัทธ์}

ให้

\[
A=\text{กรอบอ้างอิงหลัก}
\]

\[
B=\text{กรอบอ้างอิงที่เคลื่อนที่}
\]

\[
P=\text{วัตถุที่สนใจ}
\]

ตำแหน่งของ $P$ เทียบกับ $A$ เขียนได้เป็น

\[
\vec{r}_{PA}=\vec{r}_{PB}+\vec{r}_{BA}
\]

ความหมายคือ

\[
\text{ตำแหน่งของ }P\text{ เทียบกับ }A
=
\text{ตำแหน่งของ }P\text{ เทียบกับ }B
+
\text{ตำแหน่งของ }B\text{ เทียบกับ }A
\]

\subsection{ความเร็วสัมพัทธ์}

หาอนุพันธ์เทียบกับเวลา

\[
\frac{d\vec{r}_{PA}}{dt}
=
\frac{d\vec{r}_{PB}}{dt}
+
\frac{d\vec{r}_{BA}}{dt}
\]

จึงได้

\[
\vec{v}_{PA}=\vec{v}_{PB}+\vec{v}_{BA}
\]

\begin{tcolorbox}[title=สมการความเร็วสัมพัทธ์]
\[
\vec{v}_{PA}=\vec{v}_{PB}+\vec{v}_{BA}
\]
\end{tcolorbox}

\subsection{ความเร่งสัมพัทธ์}

หาอนุพันธ์ของสมการความเร็วอีกครั้ง

\[
\vec{a}_{PA}=\vec{a}_{PB}+\vec{a}_{BA}
\]

ถ้ากรอบ $B$ เคลื่อนที่ด้วยความเร็วคงตัวเทียบกับ $A$

\[
\vec{a}_{BA}=0
\]

ดังนั้น

\[
\vec{a}_{PA}=\vec{a}_{PB}
\]

\begin{tcolorbox}[title=สมการความเร่งสัมพัทธ์]
\[
\vec{a}_{PA}=\vec{a}_{PB}+\vec{a}_{BA}
\]
\end{tcolorbox}

\begin{examplebox}{ตัวอย่าง: เรือข้ามแม่น้ำ}
เรือลำหนึ่งแล่นด้วยความเร็ว $5\,\si{m/s}$ เทียบกับน้ำ ไปทางทิศเหนือ และน้ำไหลไปทางทิศตะวันออกด้วยความเร็ว $3\,\si{m/s}$ จงหาความเร็วของเรือเทียบกับฝั่ง

กำหนด

\[
\vec{v}_{\text{เรือเทียบน้ำ}}=5\hat{j}
\]

\[
\vec{v}_{\text{น้ำเทียบฝั่ง}}=3\hat{i}
\]

ดังนั้น

\[
\vec{v}_{\text{เรือเทียบฝั่ง}}
=
\vec{v}_{\text{เรือเทียบน้ำ}}
+
\vec{v}_{\text{น้ำเทียบฝั่ง}}
\]

\[
\vec{v}_{\text{เรือเทียบฝั่ง}}=3\hat{i}+5\hat{j}
\]

ขนาดคือ

\[
v=\sqrt{3^2+5^2}
\]

\[
v=\sqrt{34}\,\si{m/s}
\]
\end{examplebox}

\begin{exercisebox}{แบบฝึกหัด 4.6}
\begin{enumerate}
    \item คนเดินบนรถไฟไปข้างหน้าด้วย $1\,\si{m/s}$ เทียบกับรถไฟ และรถไฟเคลื่อนที่ไปข้างหน้าด้วย $20\,\si{m/s}$ เทียบกับพื้น จงหาความเร็วของคนเทียบกับพื้น
    \item เรือแล่นข้ามแม่น้ำไปทางเหนือด้วย $8\,\si{m/s}$ เทียบกับน้ำ แต่น้ำไหลไปทางตะวันออก $6\,\si{m/s}$ จงหาขนาดความเร็วของเรือเทียบกับฝั่ง
    \item เครื่องบินต้องการบินตรงไปทางเหนือ แต่ลมพัดไปทางตะวันออก จงอธิบายว่าต้องหันหัวเครื่องบินไปทางใด
    \item ถ้ากรอบอ้างอิงหนึ่งเคลื่อนที่ด้วยความเร็วคงตัว ความเร่งของวัตถุในสองกรอบนี้ต่างกันหรือไม่
\end{enumerate}
\end{exercisebox}

% =========================================================
\section{ชุดโจทย์รวมท้ายบท: ระดับผสมสำหรับ สอวน. ค่าย 1}

\subsection*{ระดับพื้นฐาน}

\begin{enumerate}
    \item อนุภาคมี $\vec{r}(t)=2t\hat{i}+3t^2\hat{j}$ จงหา $\vec{v}(t)$ และ $\vec{a}(t)$
    \item ยิงวัตถุด้วย $u=10\,\si{m/s}$ ทำมุม $30^\circ$ กับแนวระดับ ใช้ $g=10\,\si{m/s^2}$ จงหา $u_x$ และ $u_y$
    \item วัตถุเคลื่อนที่เป็นวงกลมรัศมี $4\,\si{m}$ ด้วยอัตราเร็ว $8\,\si{m/s}$ จงหา $a_c$
    \item เรือแล่นไปทางเหนือ $12\,\si{m/s}$ เทียบกับน้ำ และน้ำไหลไปทางตะวันออก $5\,\si{m/s}$ จงหาขนาดความเร็วเรือเทียบกับฝั่ง
\end{enumerate}

\subsection*{ระดับกลาง}

\begin{enumerate}
    \setcounter{enumi}{4}
    \item ยิงวัตถุจากพื้นด้วยความเร็ว $u$ ทำมุม $\theta$ กับแนวระดับ จงพิสูจน์ว่า
    \[
    R=\frac{u^2\sin2\theta}{g}
    \]
    เมื่อวัตถุตกที่ระดับเดิม

    \item วัตถุถูกยิงในแนวระดับจากความสูง $h$ ด้วยความเร็ว $u$ จงหาระยะทางแนวราบก่อนตกถึงพื้น

    \item รถเคลื่อนที่บนโค้งรัศมี $R$ ด้วยอัตราเร็ว $v$ และกำลังเพิ่มอัตราเร็วด้วยอัตรา $b$ จงหาขนาดความเร่งรวม

    \item เครื่องบินมีความเร็ว $u$ เทียบกับอากาศ ต้องการบินตรงไปทางเหนือ แต่ลมพัดไปทางตะวันออกด้วยความเร็ว $w$ จงหามุมที่ต้องหันหัวเครื่องบินไปทางตะวันตกของทิศเหนือ
\end{enumerate}

\subsection*{ระดับยาก}

\begin{enumerate}
    \setcounter{enumi}{8}
    \item โพรเจคไทล์ถูกยิงด้วยความเร็วต้น $u$ ทำมุม $\theta$ กับแนวระดับ และตกบนพื้นที่มีสมการ $y=-kx$ จงหาเวลาที่กระทบพื้นในรูปของ $u,\theta,k,g$

    \item วัตถุเคลื่อนที่ในระนาบโดยมี
    \[
    \vec{r}(t)=at^2\hat{i}+bt^3\hat{j}
    \]
    จงหา $\vec{v}(t)$, $\vec{a}(t)$ และขนาดของความเร็วเมื่อเวลา $t$

    \item วัตถุเคลื่อนที่บนวงกลมรัศมี $R$ โดยอัตราเร็วเพิ่มตาม
    \[
    v(t)=kt
    \]
    จงหา $a_t$, $a_r$, และขนาดความเร่งรวม ณ เวลา $t$

    \item ยิงวัตถุจากพื้นด้วยความเร็วต้น $u$ ทำมุม $\theta$ จงหาสมการเส้นทาง $y(x)$ และแสดงว่าเป็นพาราโบลา

    \item เรือแล่นข้ามแม่น้ำกว้าง $d$ ด้วยความเร็ว $u$ เทียบกับน้ำ และน้ำไหลด้วยความเร็ว $w$ ถ้าเรือเล็งตั้งฉากกับฝั่ง จงหาว่าเรือถูกพัดไปตามน้ำเป็นระยะเท่าไร
\end{enumerate}

% =========================================================
\section{เฉลยย่อท้ายบท}

\begin{enumerate}
    \item
    \[
    \vec{v}(t)=2\hat{i}+6t\hat{j}
    \]
    \[
    \vec{a}(t)=6\hat{j}
    \]

    \item
    \[
    u_x=10\cos30^\circ=5\sqrt{3}\,\si{m/s}
    \]
    \[
    u_y=10\sin30^\circ=5\,\si{m/s}
    \]

    \item
    \[
    a_c=\frac{v^2}{r}=\frac{8^2}{4}=16\,\si{m/s^2}
    \]

    \item
    \[
    v=\sqrt{12^2+5^2}=13\,\si{m/s}
    \]

    \item
    เวลาทั้งหมดเมื่อกลับสู่ระดับเดิมคือ
    \[
    T=\frac{2u\sin\theta}{g}
    \]
    พิสัยคือ
    \[
    R=u\cos\theta T
    \]
    ดังนั้น
    \[
    R=u\cos\theta\left(\frac{2u\sin\theta}{g}\right)
    \]
    \[
    R=\frac{u^2\sin2\theta}{g}
    \]

    \item
    แนวดิ่ง
    \[
    h=\frac{1}{2}gt^2
    \]
    ดังนั้น
    \[
    t=\sqrt{\frac{2h}{g}}
    \]
    แนวราบ
    \[
    x=ut
    \]
    จึงได้
    \[
    x=u\sqrt{\frac{2h}{g}}
    \]

    \item
    \[
    a_t=b
    \]
    \[
    a_r=\frac{v^2}{R}
    \]
    ดังนั้น
    \[
    a=\sqrt{b^2+\left(\frac{v^2}{R}\right)^2}
    \]

    \item
    ต้องให้ส่วนประกอบทางตะวันออกของความเร็วเครื่องบินหักล้างลม
    \[
    u\sin\alpha=w
    \]
    ดังนั้น
    \[
    \alpha=\sin^{-1}\left(\frac{w}{u}\right)
    \]

    \item
    สมการโพรเจคไทล์คือ
    \[
    x=u\cos\theta\,t
    \]
    \[
    y=u\sin\theta\,t-\frac{1}{2}gt^2
    \]
    พื้นคือ
    \[
    y=-kx
    \]
    แทน $x=u\cos\theta t$
    \[
    u\sin\theta\,t-\frac{1}{2}gt^2=-ku\cos\theta\,t
    \]
    ย้ายข้างและตัดคำตอบ $t=0$
    \[
    u(\sin\theta+k\cos\theta)-\frac{1}{2}gt=0
    \]
    ดังนั้น
    \[
    t=\frac{2u(\sin\theta+k\cos\theta)}{g}
    \]

    \item
    \[
    \vec{v}(t)=2at\hat{i}+3bt^2\hat{j}
    \]
    \[
    \vec{a}(t)=2a\hat{i}+6bt\hat{j}
    \]
    \[
    v=\sqrt{(2at)^2+(3bt^2)^2}
    \]

    \item
    \[
    a_t=\frac{dv}{dt}=k
    \]
    \[
    a_r=\frac{v^2}{R}=\frac{k^2t^2}{R}
    \]
    \[
    a=\sqrt{k^2+\left(\frac{k^2t^2}{R}\right)^2}
    \]

    \item
    จาก
    \[
    x=u\cos\theta\,t
    \]
    ได้
    \[
    t=\frac{x}{u\cos\theta}
    \]
    แทนใน
    \[
    y=u\sin\theta\,t-\frac{1}{2}gt^2
    \]
    จะได้
    \[
    y=x\tan\theta-\frac{gx^2}{2u^2\cos^2\theta}
    \]
    มีพจน์ $x^2$ จึงเป็นพาราโบลา

    \item
    เวลาข้ามแม่น้ำคือ
    \[
    t=\frac{d}{u}
    \]
    ระยะที่ถูกพัดตามน้ำคือ
    \[
    x=wt
    \]
    ดังนั้น
    \[
    x=\frac{wd}{u}
    \]
\end{enumerate}

% =========================================================
\section{สรุปสูตรสำคัญ}

\begin{tcolorbox}[title=สรุปบทที่ 4: การเคลื่อนที่ในสองมิติ]
เวกเตอร์ตำแหน่ง ความเร็ว และความเร่ง

\[
\vec{r}(t)=x(t)\hat{i}+y(t)\hat{j}
\]

\[
\vec{v}(t)=\frac{d\vec{r}}{dt}
=
\frac{dx}{dt}\hat{i}
+
\frac{dy}{dt}\hat{j}
\]

\[
\vec{a}(t)=\frac{d\vec{v}}{dt}
=
\frac{d^2\vec{r}}{dt^2}
\]

เมื่อความเร่งคงตัว

\[
\vec{v}_f=\vec{v}_i+\vec{a}t
\]

\[
\vec{r}_f=\vec{r}_i+\vec{v}_it+\frac{1}{2}\vec{a}t^2
\]

โพรเจคไทล์

\[
u_x=u\cos\theta
\]

\[
u_y=u\sin\theta
\]

\[
x=u\cos\theta\,t
\]

\[
y=u\sin\theta\,t-\frac{1}{2}gt^2
\]

\[
v_x=u\cos\theta
\]

\[
v_y=u\sin\theta-gt
\]
\end{tcolorbox}

\begin{tcolorbox}[title=สรุปบทที่ 4: การเคลื่อนที่ในสองมิติ ต่อ]
กรณียิงและตกที่ระดับเดิม

\[
t_{\text{top}}=\frac{u\sin\theta}{g}
\]

\[
H=\frac{u^2\sin^2\theta}{2g}
\]

\[
T=\frac{2u\sin\theta}{g}
\]

\[
R=\frac{u^2\sin2\theta}{g}
\]

วงกลมสม่ำเสมอ

\[
a_c=\frac{v^2}{r}
\]

\[
v=\frac{2\pi r}{T}
\]

\[
\omega=\frac{2\pi}{T}
\]

\[
v=\omega r
\]

\[
a_c=\omega^2r
\]

ความเร่งแนวสัมผัสและแนวรัศมี

\[
a_t=\left|\frac{dv}{dt}\right|
\]

\[
a_r=\frac{v^2}{r}
\]

\[
a=\sqrt{a_t^2+a_r^2}
\]

ความเร็วสัมพัทธ์

\[
\vec{r}_{PA}=\vec{r}_{PB}+\vec{r}_{BA}
\]

\[
\vec{v}_{PA}=\vec{v}_{PB}+\vec{v}_{BA}
\]

\[
\vec{a}_{PA}=\vec{a}_{PB}+\vec{a}_{BA}
\]
\end{tcolorbox}

% =========================================================
% =========================================================
\section{ข้อสอบจริง สอวน. ที่ใช้ความรู้บทการเคลื่อนที่ในสองมิติ}


% =========================================================
\subsection{ปี 2567 ข้อ 2: โพรเจคไทล์ตกบนพื้นเอียง}

\vspace{1.5em}

\IfFileExists{img/posn67q2.png}{
\begin{center}
    \includegraphics[width=1\linewidth]{img/posn67q2.png}
\end{center}
}{
\begin{tcolorbox}[title=ตำแหน่งภาพที่แนะนำ]
ให้ตัดภาพข้อสอบปี 2567 ข้อ 2 แล้วบันทึกเป็น

\[
\texttt{img/posn67q2.png}
\]
\end{tcolorbox}
}

\vspace{1.5em}

\begin{examplebox}{เฉลยละเอียด}
โจทย์นี้เป็นโพรเจคไทล์ที่ตกบนพื้นเอียง ดังนั้นห้ามใช้สูตรพิสัย

\[
R=\frac{u^2\sin2\theta}{g}
\]

เพราะสูตรนี้ใช้ได้เฉพาะกรณีที่จุดเริ่มต้นและจุดตกอยู่ระดับเดียวกันเท่านั้น

จากรูป วัตถุถูกยิงด้วยความเร็วต้น $u$ ทำมุม

\[
30^\circ
\]

กับแนวระดับ และพื้นเอียงทำมุม

\[
30^\circ
\]

ใต้แนวระดับ

เลือกแกน $x$ แนวราบไปทางขวา และแกน $y$ แนวดิ่งขึ้น

แตกความเร็วต้นเป็นองค์ประกอบ

\[
u_x=u\cos30^\circ
\]

\[
u_y=u\sin30^\circ
\]

ดังนั้นสมการโพรเจคไทล์คือ

\[
x=u\cos30^\circ t
\]

\[
y=u\sin30^\circ t-\frac{1}{2}gt^2
\]

พื้นเอียงทำมุม $30^\circ$ ใต้แนวระดับ จึงมีสมการ

\[
y=-x\tan30^\circ
\]

แทน

\[
x=u\cos30^\circ t
\]

ลงในสมการพื้นเอียง

\[
y=-(u\cos30^\circ t)\tan30^\circ
\]

เพราะ

\[
\cos30^\circ\tan30^\circ=\sin30^\circ
\]

จึงได้

\[
y=-u\sin30^\circ t
\]

ที่จุดกระทบ สมการโพรเจคไทล์และสมการพื้นเอียงต้องให้ค่า $y$ เท่ากัน

\[
u\sin30^\circ t-\frac{1}{2}gt^2=-u\sin30^\circ t
\]

ย้ายข้าง

\[
2u\sin30^\circ t-\frac{1}{2}gt^2=0
\]

ใช้

\[
\sin30^\circ=\frac{1}{2}
\]

จะได้

\[
ut-\frac{1}{2}gt^2=0
\]

ไม่เอาคำตอบ $t=0$ เพราะเป็นจุดเริ่มต้น

\[
t=\frac{2u}{g}
\]

หาระยะตามแนวราบ

\[
x=u\cos30^\circ\left(\frac{2u}{g}\right)
\]

\[
x=\frac{2u^2\cos30^\circ}{g}
\]

\[
x=\frac{\sqrt{3}u^2}{g}
\]

ให้ $OP$ เป็นระยะตามพื้นเอียง

\[
x=OP\cos30^\circ
\]

ดังนั้น

\[
OP=\frac{x}{\cos30^\circ}
\]

\[
OP=
\frac{\frac{\sqrt{3}u^2}{g}}{\frac{\sqrt{3}}{2}}
\]

\[
\boxed{OP=\frac{2u^2}{g}}
\]

ต่อไปหามุมที่ความเร็วกระทบทำกับแนวดิ่ง

ความเร็วแนวราบคือ

\[
v_x=u\cos30^\circ=\frac{\sqrt{3}}{2}u
\]

ความเร็วแนวดิ่งคือ

\[
v_y=u\sin30^\circ-gt
\]

แทน

\[
t=\frac{2u}{g}
\]

จะได้

\[
v_y=\frac{1}{2}u-g\left(\frac{2u}{g}\right)
\]

\[
v_y=-\frac{3}{2}u
\]

ให้ $\beta$ เป็นมุมที่ความเร็วทำกับแนวระดับลงด้านล่าง

\[
\tan\beta=\frac{|v_y|}{v_x}
\]

\[
\tan\beta=
\frac{\frac{3}{2}u}{\frac{\sqrt{3}}{2}u}
\]

\[
\tan\beta=\sqrt{3}
\]

ดังนั้น

\[
\beta=60^\circ
\]

แนวดิ่งลงทำมุม $90^\circ$ กับแนวระดับ ดังนั้นมุมระหว่างความเร็วกับแนวดิ่งลงคือ

\[
90^\circ-60^\circ=30^\circ
\]

จึงได้

\[
\boxed{\phi=30^\circ}
\]
\end{examplebox}

% =========================================================
\subsection{ปี 2568 ข้อ 1: โพรเจคไทล์และทิศของความเร็วที่จุดหนึ่ง}

\vspace{1.5em}

\IfFileExists{img/posn68q1.png}{
\begin{center}
    \includegraphics[width=1\linewidth]{img/posn68q1.png}
\end{center}
}{
\begin{tcolorbox}[title=ตำแหน่งภาพที่แนะนำ]
ให้ตัดภาพข้อสอบปี 2568 ข้อ 1 แล้วบันทึกเป็น

\[
\texttt{img/posn68q1.png}
\]
\end{tcolorbox}
}

\vspace{1.5em}

\begin{examplebox}{เฉลยละเอียด}
จากรูป โพรเจคไทล์ออกจากจุด $O$ ด้วยความเร็วต้น $u$ ทำมุม $45^\circ$ กับแนวดิ่ง ดังนั้นทำมุม $45^\circ$ กับแนวระดับเช่นกัน

แตกความเร็วต้นเป็นองค์ประกอบ

\[
u_x=u\cos45^\circ=\frac{u}{\sqrt{2}}
\]

\[
u_y=u\sin45^\circ=\frac{u}{\sqrt{2}}
\]

ในการเคลื่อนที่แบบโพรเจคไทล์โดยไม่คิดแรงต้านอากาศ ความเร็วแนวราบคงตัว

\[
v_x=u_x=\frac{u}{\sqrt{2}}
\]

ที่จุด $A$ จากรูป ทิศของความเร็วทำมุม $60^\circ$ กับแนวดิ่ง ดังนั้นทำมุม

\[
30^\circ
\]

กับแนวระดับ

จึงมี

\[
\tan30^\circ=\frac{v_y}{v_x}
\]

ดังนั้น

\[
v_y=v_x\tan30^\circ
\]

แทนค่า

\[
v_y=\frac{u}{\sqrt{2}}\left(\frac{1}{\sqrt{3}}\right)
\]

\[
v_y=\frac{u}{\sqrt{6}}
\]

ต้องการหาความสูงของจุด $A$ จากระดับเริ่มต้น ใช้สมการแนวดิ่งที่ไม่ต้องใช้เวลา

\[
v_y^2=u_y^2-2gy
\]

จัดรูป

\[
y=\frac{u_y^2-v_y^2}{2g}
\]

โดย

\[
u_y^2=\left(\frac{u}{\sqrt{2}}\right)^2=\frac{u^2}{2}
\]

และ

\[
v_y^2=\left(\frac{u}{\sqrt{6}}\right)^2=\frac{u^2}{6}
\]

แทนค่า

\[
y=\frac{\frac{u^2}{2}-\frac{u^2}{6}}{2g}
\]

\[
y=\frac{\frac{u^2}{3}}{2g}
\]

ดังนั้น

\[
\boxed{y=\frac{u^2}{6g}}
\]
\end{examplebox}

% =========================================================
\subsection{ปี 2564 ข้อ 30: โพรเจคไทล์เฉียดยอดกำแพง}

\vspace{1.5em}

\IfFileExists{img/posn64q30.png}{
\begin{center}
    \includegraphics[width=1\linewidth]{img/posn64q30.png}
\end{center}
}{
\begin{tcolorbox}[title=ตำแหน่งภาพที่แนะนำ]
ให้ตัดภาพข้อสอบปี 2564 ข้อ 30 แล้วบันทึกเป็น

\[
\texttt{img/posn64q30.png}
\]
\end{tcolorbox}
}

\vspace{1.5em}

\begin{examplebox}{เฉลยละเอียด}
จากรูป โพรเจคไทล์ถูกยิงจากจุด $A$ ด้วยความเร็วต้น $u$ ทำมุม $\theta$ กับพื้นระดับ และเฉียดจุดยอดกำแพงที่จุด $B$

กำแพงอยู่ห่างจากจุดยิงในแนวราบเป็นระยะ

\[
L
\]

และจุดยอดกำแพงสูงจากพื้นเป็นระยะ

\[
H
\]

จากรูป จุด $B$ เป็นจุดสูงสุดของวิถีพอดี ดังนั้นที่จุด $B$

\[
v_y=0
\]

แตกความเร็วต้นเป็นองค์ประกอบ

\[
u_x=u\cos\theta
\]

\[
u_y=u\sin\theta
\]

เวลาที่วัตถุขึ้นถึงจุดสูงสุดหาได้จากแนวดิ่ง

\[
v_y=u_y-gt
\]

ที่จุดสูงสุด

\[
0=u\sin\theta-gt
\]

ดังนั้น

\[
t=\frac{u\sin\theta}{g}
\]

พิจารณาแนวราบ วัตถุเคลื่อนที่ถึงกำแพงที่ระยะ $L$

\[
L=u_x t
\]

แทนค่า

\[
L=u\cos\theta\left(\frac{u\sin\theta}{g}\right)
\]

ดังนั้น

\[
L=\frac{u^2\sin\theta\cos\theta}{g}
\]

พิจารณาแนวดิ่ง ความสูงสูงสุดคือ

\[
H=\frac{u_y^2}{2g}
\]

แทนค่า $u_y=u\sin\theta$

\[
H=\frac{u^2\sin^2\theta}{2g}
\]

นำสมการของ $H$ หารด้วยสมการของ $L$

\[
\frac{H}{L}
=
\frac{\frac{u^2\sin^2\theta}{2g}}
{\frac{u^2\sin\theta\cos\theta}{g}}
\]

ตัดตัวร่วม

\[
\frac{H}{L}
=
\frac{1}{2}\frac{\sin\theta}{\cos\theta}
\]

ดังนั้น

\[
\frac{H}{L}=\frac{1}{2}\tan\theta
\]

จึงได้

\[
\boxed{\tan\theta=\frac{2H}{L}}
\]

ต่อไปหา $u^2$

จาก

\[
H=\frac{u^2\sin^2\theta}{2g}
\]

จัดรูป

\[
u^2=\frac{2gH}{\sin^2\theta}
\]

ใช้ความสัมพันธ์

\[
\sin^2\theta=\frac{\tan^2\theta}{1+\tan^2\theta}
\]

และ

\[
\tan\theta=\frac{2H}{L}
\]

ดังนั้น

\[
\sin^2\theta
=
\frac{\left(\frac{2H}{L}\right)^2}
{1+\left(\frac{2H}{L}\right)^2}
\]

\[
\sin^2\theta
=
\frac{4H^2}{L^2+4H^2}
\]

แทนกลับในสมการของ $u^2$

\[
u^2
=
\frac{2gH}{\frac{4H^2}{L^2+4H^2}}
\]

\[
u^2
=
2gH\left(\frac{L^2+4H^2}{4H^2}\right)
\]

ดังนั้น

\[
\boxed{u^2=\frac{g(L^2+4H^2)}{2H}}
\]
\end{examplebox}

% =========================================================
\subsection{ปี 2565 ข้อ 27: การเล็งสูงกว่าจุดเป้าหมายเล็กน้อย}

\vspace{1.5em}

\IfFileExists{img/posn65q27.png}{
\begin{center}
    \includegraphics[width=1\linewidth]{img/posn65q27.png}
\end{center}
}{
\begin{tcolorbox}[title=ตำแหน่งภาพที่แนะนำ]
ให้ตัดภาพข้อสอบปี 2565 ข้อ 27 แล้วบันทึกเป็น

\[
\texttt{img/posn65q27.png}
\]
\end{tcolorbox}
}

\vspace{1.5em}

\begin{examplebox}{เฉลยละเอียด}
กำหนดให้อัตราเร็วต้นของโพรเจคไทล์เท่ากับ

\[
v_0
\]

ต้องการให้โพรเจคไทล์ชนจุด $B$ บนกำแพง ซึ่งอยู่ห่างจากจุดยิงในแนวราบเป็นระยะ

\[
D
\]

ถ้าไม่มีแรงโน้มถ่วง เราสามารถเล็งตรงไปที่จุด $B$ ได้พอดี แต่เมื่อมีแรงโน้มถ่วง วัตถุจะตกต่ำลงระหว่างทาง ดังนั้นต้องเล็งสูงกว่าจุด $B$ เป็นมุมเล็ก ๆ

\[
\delta\theta
\]

ให้เส้นตรงจากจุดยิงไปยังจุด $B$ ทำมุม $\theta$ กับแนวระดับ และให้มุมยิงจริงเป็น

\[
\theta+\delta\theta
\]

สมการโพรเจคไทล์ที่ตำแหน่งแนวราบ $x=D$ คือ

\[
y(D)
=
D\tan(\theta+\delta\theta)
-
\frac{gD^2}{2v_0^2\cos^2(\theta+\delta\theta)}
\]

เนื่องจาก $\delta\theta$ มีค่าน้อยมาก ใช้การประมาณอันดับหนึ่ง

\[
\tan(\theta+\delta\theta)
\approx
\tan\theta+\delta\theta\sec^2\theta
\]

และในพจน์แรงโน้มถ่วง ใช้

\[
\cos^2(\theta+\delta\theta)\approx \cos^2\theta
\]

ดังนั้น

\[
y(D)
\approx
D\tan\theta
+
D\delta\theta\sec^2\theta
-
\frac{gD^2}{2v_0^2\cos^2\theta}
\]

แต่จุด $B$ อยู่บนเส้นเล็งเดิม จึงมีความสูง

\[
y_B=D\tan\theta
\]

เพื่อให้วัตถุชนจุด $B$ ต้องมี

\[
y(D)=y_B
\]

ดังนั้น

\[
D\tan\theta
+
D\delta\theta\sec^2\theta
-
\frac{gD^2}{2v_0^2\cos^2\theta}
=
D\tan\theta
\]

ตัด $D\tan\theta$ ทั้งสองข้าง

\[
D\delta\theta\sec^2\theta
=
\frac{gD^2}{2v_0^2\cos^2\theta}
\]

เพราะ

\[
\sec^2\theta=\frac{1}{\cos^2\theta}
\]

จึงได้

\[
\frac{D\delta\theta}{\cos^2\theta}
=
\frac{gD^2}{2v_0^2\cos^2\theta}
\]

ตัด $\cos^2\theta$ และ $D$

\[
\delta\theta=\frac{gD}{2v_0^2}
\]

ดังนั้นมุมที่ต้องเล็งสูงกว่าจุด $B$ โดยประมาณคือ

\[
\boxed{\delta\theta\approx \frac{gD}{2v_0^2}}
\]

โดยมุมนี้มีหน่วยเป็นเรเดียน
\end{examplebox}

% =========================================================
\subsection{ปี 2560 ข้อ 8: ผลของแรงต้านอากาศต่อวิถีโพรเจคไทล์}

\vspace{1.5em}

\IfFileExists{img/posn60q8.png}{
\begin{center}
    \includegraphics[width=1\linewidth]{img/posn60q8.png}
\end{center}
}{
\begin{tcolorbox}[title=ตำแหน่งภาพที่แนะนำ]
ให้ตัดภาพข้อสอบปี 2560 ข้อ 8 แล้วบันทึกเป็น

\[
\texttt{img/posn60q8.png}
\]

ข้อนี้เป็นข้อเสริม เพราะมีแรงต้านอากาศ ซึ่งเกินแบบจำลองโพรเจคไทล์พื้นฐานเล็กน้อย
\end{tcolorbox}
}

\vspace{1.5em}

\begin{examplebox}{เฉลยแนวคิด}
ในบทโพรเจคไทล์พื้นฐาน เราถือว่าไม่มีแรงต้านอากาศ ดังนั้น

\[
a_x=0
\]

\[
a_y=-g
\]

และวิถีเป็นพาราโบลา

แต่ถ้ามีแรงต้านอากาศ แรงต้านจะมีทิศตรงข้ามกับทิศการเคลื่อนที่เสมอ ทำให้เกิดผลสำคัญดังนี้

\[
\text{ความเร็วแนวราบลดลง}
\]

\[
\text{วัตถุไปได้ไม่ไกลเท่ากรณีไม่มีแรงต้าน}
\]

\[
\text{ความสูงสูงสุดต่ำลง}
\]

\[
\text{ขาลงมักชันกว่าและไม่สมมาตรกับขาขึ้น}
\]

ดังนั้นกราฟที่ถูกต้องต้องอยู่ต่ำกว่าเส้นประซึ่งแทนกรณีไม่มีแรงต้าน และตกถึงพื้นก่อนหรือมีพิสัยสั้นกว่า

\[
\boxed{
\text{เลือกเส้นที่อยู่ต่ำกว่าเส้นประ และมีพิสัยสั้นกว่ากรณีไม่มีแรงต้าน}
}
\]


\end{examplebox}
